\documentclass[../../main.tex]{subfiles}
\begin{document}

\begin{flushright}
  \footnotesize\textit{【自動文字起こし・要確認】}
\end{flushright}

[解]
(i) $t > 0$から, AM-GM
\begin{align}
  f(t) \ge 2\sqrt{t \cdot \dfrac{1}{t}} + \sqrt{2\left(t \cdot \dfrac{1}{t} + 1\right)} = 2 + \sqrt{3}
\end{align}
(等号成立は$t=1$)
\begin{align}
  f(t)g(t) = \left(t+\dfrac{1}{t}+2\right) - \left(t+\dfrac{1}{t}+1\right) = 1
\end{align}
\begin{align}
  \therefore g(t) = \dfrac{1}{f(t)} \quad (\because f(t) \ge 2+\sqrt{3})
\end{align}
だから, (i)から, $g(t)$は$f(t) = 2+\sqrt{3}$の時, $\max \dfrac{2-\sqrt{3}}{1}$をとる. 終

(ii) $x, y > 0$から

「$a, b, c$を$3$辺とする三角形が常に存在」

$\iff$ 「$\dfrac{a}{\sqrt{xy}}, \dfrac{b}{\sqrt{xy}}, \dfrac{c}{\sqrt{xy}} \quad \text{〃}$」

である. ここで, $\dfrac{a}{\sqrt{xy}} = A$などとおき, $t = \dfrac{x}{y} \ (t>0)$とおくと
\begin{align}
  A = \sqrt{t+\dfrac{1}{t}+1}, \ B = p, \ C = \sqrt{t} + \dfrac{1}{\sqrt{t}}
\end{align}
まず, (i)から
\begin{equation}
  \max(C-A) = 2-\sqrt{3}, \ \min(A+C) = 2+\sqrt{3} \label{eq:1}
\end{equation}
だから, もとめる条件は\eqref{eq:1}が成立するような$A, C$の時での成立が必要で,
\begin{align}
  2-\sqrt{3} < p < 2+\sqrt{3}
\end{align}
逆にこの時, $A+C \ge 2+\sqrt{3}, \ C-A \le 2-\sqrt{3}$だから$0 < C-A < p < A+C$となり, たしかに$A, B, C$を3辺とする三角形をつくることが出来, ($\because$ 三角不等式) 十分である.

以上から
\begin{align}
  2-\sqrt{3} < p < 2+\sqrt{3}
\end{align}

\end{document}
